\documentclass[12pt,a4paper]{article}

\usepackage[utf8]{inputenc}
\usepackage[english]{babel}
\usepackage{amsmath,amssymb}
\usepackage{graphicx}
\usepackage{booktabs}
\usepackage{hyperref}
\usepackage{natbib}
\usepackage{float}
\usepackage{multirow}
\usepackage{caption}
\usepackage{subcaption}
\usepackage{colortbl}
\usepackage{xcolor}
\usepackage[margin=2.5cm]{geometry}
\usepackage{microtype}

\graphicspath{{../../slides_export/}}

\definecolor{bestrow}{RGB}{200,230,200}
\definecolor{greyrow}{RGB}{242,242,242}

\title{Biotic Interaction Detection from Scientific Literature:\\
Knowledge Distillation and Multi-task Learning\\
with Named Entity Recognition}

\author{
Esteban Gaillac\textsuperscript{1,*} \\[4pt]
\textsuperscript{1}Haute \'Ecole de Gestion, HES-SO Geneva \\
\textsuperscript{*}Correspondence: esteban.gaillac@hesge.ch
}

\date{June 2026}

\begin{document}

\maketitle

% ─────────────────────────────────────────────────────────────────────────────
\begin{abstract}
Detecting biotic interactions from scientific text is a prerequisite for
automated knowledge graph construction from biodiversity literature.
We present a multi-task transformer classifier that achieves F1\,=\,0.868
(P\,=\,0.793, R\,=\,0.958) on a 100-sentence ecological test set,
outperforming both a fine-tuned BiomedBERT baseline (F1\,=\,0.788) and
a BiomedBERT\,$\times$\,FLAN-T5-base ensemble (F1\,=\,0.865) at 3.3$\times$
lower inference cost.

The model is built on three components. First, a two-stage teacher chain
in which a local 122-billion-parameter model (Qwen3.5-122B) annotates
44,178 harvested sentences with binary labels, and a trained ensemble
subsequently generates calibrated soft probability labels for knowledge
distillation. Second, a multi-task BiomedBERT model that jointly classifies
sentences and annotates tokens with biological entity roles (HOST, PATHOGEN,
SPECIES, INT) using automatically derived BIO supervision, with a two-phase
training strategy that primes the shared encoder via NER before introducing
the classification objective. Third, we show that soft labels are the single
largest performance lever in the pipeline: switching from hard to soft
cross-entropy on identical architecture and data yields a +0.052 F1 gain.

A distribution analysis reveals that 77\% of test-set positives belong to
interaction types underrepresented in the PubMed-sourced training corpus,
explaining the current performance ceiling at F1\,$\approx$\,0.87 and
providing a clear target for future data collection. The classifier is
deployed upstream of BiotXplorer
(\url{https://biotxplorer.sibils.org/}), the SIBiLS \citep{mottin2021sibils}
semantic search platform for biotic interactions, and is designed as the
first component of a downstream knowledge graph construction pipeline.
\end{abstract}

\textbf{Keywords:} biotic interactions, knowledge distillation, multi-task
learning, named entity recognition, BiomedBERT, knowledge graph,
ecological text mining, GloBI

% ─────────────────────────────────────────────────────────────────────────────
\section{Introduction}
\label{sec:intro}

Species interactions shape ecosystem structure, drive disease dynamics,
and determine the response of ecological communities to perturbation.
The Global Biotic Interactions database (GloBI; \citealt{poelen2014global})
aggregates over 13 million structured interaction records from more than 700
data sources, yet the vast majority of interaction knowledge remains encoded
in unstructured scientific literature distributed across millions of
publications. Biodiversity PMC, the SIBiLS \citep{mottin2021sibils} literature
corpus covering PubMed Central content relevant to biodiversity, alone indexes
more than 40 million articles. Systematic curation of this corpus is
infeasible at human scale.

Sentence-level classification of biotic interactions is the entry point
for automated literature mining. Given a sentence, the classifier determines
whether it describes a biotic interaction between two named organisms:
predation, parasitism, pathogen infection, pollination, herbivory, mutualism,
or symbiosis. This binary decision feeds two downstream applications.

\paragraph{Upstream: BiotXplorer.}
BiotXplorer \citep{sibils2023biotxplorer} is a semantic search engine for
biotic interactions in the biomedical literature built on the SIBiLS
infrastructure \citep{mottin2021sibils}. The classifier serves two complementary
roles in this context. First, it acts as an upstream filter, screening millions
of candidate sentences before indexing. Second, and more precisely, it removes
false positives from co-occurrence-based candidate selection: given a triplet
(species\,1, species\,2, interaction term) detected in a sentence, the
classifier decides whether the sentence actually describes the interaction or
merely co-mentions the terms in a non-interacting context. Both use cases
reduce index noise while preserving comprehensive coverage.

\paragraph{Downstream: knowledge graph construction.}
Beyond search, structured interaction triples extracted from accepted sentences
can populate biodiversity knowledge graphs \citep{chen2023biodiversity}
supporting species-level reasoning about disease spread, invasion risk,
food web dynamics, and conservation planning \citep{poisot2021global}.
The classifier is the first component in such a pipeline; downstream relation
extraction maps accepted sentences to typed triples of the form
(subject entity, relation, object entity, confidence, provenance).

Previous NLP work in the biomedical domain has focused primarily on
protein-protein interactions \citep{bunescu2005comparative} and drug-gene
relationships \citep{lim2018drug}, with ecological interaction types
(predation, parasitism, pollination, herbivory) receiving comparatively less
attention. The ecological literature presents additional challenges: it spans
thousands of journals, uses varied terminology, and describes field
observations that differ structurally from the experimental prose dominant
in PubMed \citep{badenes2020detecting}.

We address this gap with three primary contributions:
\begin{enumerate}
    \item A two-stage teacher chain combining large-LLM binary annotation
    with ensemble-based soft probability distillation, yielding a
    44K-sentence training corpus with calibrated labels.
    \item A multi-task BiomedBERT architecture with automatic NER auxiliary
    supervision, trained in two phases to prime the shared encoder.
    \item A quantitative distribution mismatch analysis identifying the
    interaction types responsible for the current performance ceiling and
    guiding future data collection.
\end{enumerate}

% ─────────────────────────────────────────────────────────────────────────────
\section{Methods}
\label{sec:methods}

\subsection{Task Definition}

Given a sentence $s$ from scientific literature, predict $y \in \{0, 1\}$
where $y=1$ indicates that $s$ describes a biotic interaction between two
or more named organisms. Interaction types in scope: predation, parasitism,
pathogen infection, pollination, herbivory, mutualism, symbiosis, seed
dispersal, competition, and vector-borne transmission.

\subsection{Training Corpus Construction}
\label{sec:data}

\paragraph{Sentence sources.}
Two sentence origins are combined into a single 44,178-sentence corpus.

\emph{GloBI-derived templates (v1--v7, synthetic).}
Structured triples from GloBI (e.g.\ \textit{parasitizes(Anisakis simplex,
Gadus morhua)}) are injected into hand-written sentence templates covering
nine interaction categories. In the final template version (v7), every
positive sentence was individually validated by the Claude API (Anthropic)
to remove formulaic or linguistically implausible candidates, yielding
25,081 clean sentences (7,251 positive; F1\,=\,0.788).

\emph{Real PMC sentences (v10+).}
Articles co-mentioning known GloBI species pairs were queried from Biodiversity
PMC \citep{mottin2021sibils} ($\approx$39K sentences in total, queried both
by species co-occurrence and by interaction category via the SIBiLS API).
Co-occurrence sentences were extracted and labeled by a local Qwen3.5-122B
model using the prompt:
\textit{``Does this sentence describe a biotic interaction between two named
organisms? Answer yes or no only.''}
Of 44,178 harvested sentences, 4,065 (9.2\%) were accepted as positives.

\paragraph{Soft label generation.}
The accepted 44,178 sentences were re-scored by the trained BiomedBERT
$\times$ FLAN-T5-base ensemble, producing a soft probability
$p_\text{ensemble} \in [0,1]$ per sentence. These soft labels constitute
the distillation teacher signal (Section~\ref{sec:distillation}).
Figure~\ref{fig:version_f1} summarises F1 on the test set across dataset
configurations.

\begin{figure}[H]
    \centering
    \includegraphics[width=0.85\textwidth]{fig_01_version_f1}
    \caption{Test F1 by training data configuration.
    The best result (rightmost) is the 44K Qwen-labeled corpus with
    ensemble soft labels.}
    \label{fig:version_f1}
\end{figure}

\subsection{Model Architecture Evolution}
\label{sec:architectures}

We evaluated five approaches in order of increasing complexity.

\paragraph{Classical baselines.}
SVM and Random Forest on TF-IDF features achieved F1\,$\approx$\,0.62,
establishing the need for contextual representations.

\paragraph{Fine-tuned BiomedBERT.}
We use the BiomedBERT model \citep{gu2021domain}\\
(\texttt{microsoft/BiomedNLP-BiomedBERT-base-uncased-abstract-fulltext}),
a BERT-base model (110M parameters, 12 transformer
layers, 768-dimensional hidden states) pre-trained from scratch on PubMed
abstracts and PubMed Central full-text articles. A linear classification
head on the [CLS] token representation is fine-tuned using 5-fold
cross-validation, AdamW optimiser ($\eta = 2 \times 10^{-5}$, weight
decay\,=\,0.01), and threshold optimisation on the validation split.
BiomedBERT trained on v7 templates achieves F1\,=\,0.788.

\paragraph{BiomedBERT\,$\times$\,FLAN-T5-base ensemble.}
FLAN-T5-base \citep{chung2022scaling}, a sequence-to-sequence model trained
with instruction tuning, was fine-tuned as a binary classifier by prompting
for a yes/no answer and scoring the first generated token. The geometric
mean of BiomedBERT and FLAN-T5-base predicted probabilities, with joint
threshold optimisation, achieves F1\,=\,0.865. The gain is attributable to
model complementarity: BiomedBERT handles precise, template-like language
while FLAN-T5-base generalises better to paraphrased ecological sentences.

\paragraph{Knowledge distillation.}
Section~\ref{sec:distillation}.

\paragraph{Multi-task BiomedBERT with NER.}
Section~\ref{sec:multitask}.

Figure~\ref{fig:model_evolution} summarises the F1 progression across architectures.

\begin{figure}[H]
    \centering
    \includegraphics[width=0.85\textwidth]{fig_02_model_evolution}
    \caption{Test F1 across model architecture stages. Annotated arrows indicate
    the primary driver of each improvement.}
    \label{fig:model_evolution}
\end{figure}

\subsection{Knowledge Distillation}
\label{sec:distillation}

Knowledge distillation \citep{hinton2015distilling} transfers calibrated
uncertainty from a teacher to a student by training on soft probability
distributions rather than hard labels \citep{bucilua2006model}.
A teacher score of 0.7 instructs the student that a sentence is
\emph{probably} positive but not certain; the student is penalised for
predicting 0.4, not merely for predicting the wrong class. This
information transfer is especially valuable for borderline sentences.

\paragraph{Training objective.}
The student BiomedBERT is trained using KL divergence against the ensemble
teacher:
\begin{equation}
    \mathcal{L}_\text{KL} = T^2 \cdot \mathrm{KL}\!\left(
        \sigma(\mathbf{z}_s / T) \;\|\; \sigma(\mathbf{z}_t / T)
    \right)
\end{equation}
where $T$ is the temperature, $\mathbf{z}_s$ and $\mathbf{z}_t$ are student
and teacher logits, and $\sigma$ denotes the softmax function.
Rows without soft labels use hard cross-entropy. The combined loss is
$\mathcal{L} = \alpha \cdot \mathcal{L}_\text{KL} + (1{-}\alpha) \cdot \mathcal{L}_\text{CE}$.

Optimal hyperparameters identified by search: $T{=}2$, $\alpha{=}0.5$
(Table~\ref{tab:distillation}). Temperature $T{=}2$ spreads the teacher
distribution sufficiently to transfer calibrated uncertainty without
over-smoothing; $T{=}4$ over-smooths while $T{=}1.5$ loses the calibration
benefit. The student (distilled\_v2) achieves F1\,=\,0.808, a gain of +0.020
over the v7 fine-tuned baseline.

\begin{table}[H]
\centering
\caption{Knowledge distillation hyperparameter search.
All variants use BiomedBERT as student except v4 (DistilBERT) and v5 (SciBERT).}
\label{tab:distillation}
\begin{tabular}{lccc}
\toprule
\textbf{Version} & $T$ & $\alpha$ & \textbf{Test F1} \\
\midrule
v1 & 4 & 0.7 & 0.785 \\
\rowcolor{bestrow}
\textbf{v2 (best)} & \textbf{2} & \textbf{0.5} & \textbf{0.808} \\
v3 & 4 & 0.9 & 0.808 \\
v4 (DistilBERT) & 2 & 0.5 & 0.792 \\
v5 (SciBERT) & 2 & 0.5 & 0.793 \\
v6 ($T=1.5$) & 1.5 & 0.5 & 0.779 \\
\bottomrule
\end{tabular}
\end{table}

% ─────────────────────────────────────────────────────────────────────────────
% FROM HERE: SECTION 2.5 ONWARDS — STRUCTURE ONLY
% ─────────────────────────────────────────────────────────────────────────────

\subsection{Multi-task BiomedBERT with NER Auxiliary Task}
\label{sec:multitask}

\paragraph{Architecture.}

\paragraph{NER label scheme.}

\paragraph{Two-phase training.}

\subsection{Evaluation Protocol}
\label{sec:eval}

% ─────────────────────────────────────────────────────────────────────────────
\section{Results}
\label{sec:results}

\subsection{Multi-task Ablation}

\begin{table}[H]
\centering
\small
\caption{Multi-task ablation results on the test set (100 sentences, 48 positives).
NER pt = NER pre-train epochs before joint training.
Baseline comparisons are shown below the separator.}
\label{tab:ablation}
\begin{tabular}{llccrrr}
\toprule
\textbf{Config} & \textbf{NER scheme} & $\alpha$ & \textbf{NER pt} &
\textbf{F1} & \textbf{P} & \textbf{R} \\
\midrule
basic\_a05           & basic       & 0.5 & 0 & 0.841 & 0.763 & 0.938 \\
typed\_a05           & typed       & 0.5 & 0 & 0.852 & 0.767 & 0.958 \\
full\_a03            & full        & 0.3 & 0 & 0.854 & 0.854 & 0.854 \\
full\_typed\_a05     & full\_typed & 0.5 & 0 & 0.847 & 0.746 & 0.979 \\
full\_a05\_ner2      & full        & 0.5 & 2 & 0.824 & 0.778 & 0.875 \\
full\_typed\_a03\_ner2 & full\_typed & 0.3 & 2 & 0.830 & 0.759 & 0.917 \\
\rowcolor{bestrow}
\textbf{full\_typed\_a05\_ner2} & full\_typed & 0.5 & 2 & \textbf{0.868} & \textbf{0.793} & \textbf{0.958} \\
full\_typed (5 joint ep) & full\_typed & 0.5 & 2+5 & 0.816 & 0.764 & 0.875 \\
\midrule
\rowcolor{greyrow}
v12\_hardce          & full\_typed & 0.5 & 2 & 0.816 & 0.764 & 0.875 \\
\rowcolor{greyrow}
distilled\_v2        & ---         & --- & --- & 0.808 & 0.750 & 0.875 \\
\rowcolor{greyrow}
Ensemble BERT$\times$T5 & ---      & --- & --- & 0.865 & ---   & ---   \\
\bottomrule
\end{tabular}
\end{table}

\paragraph{1. Soft labels are the largest single lever (+0.052 F1).}

\paragraph{2. NER pre-training is necessary.}

\paragraph{3. Richer NER schemes perform better.}

\subsection{Statistical Significance}

\begin{figure}[H]
    \centering
    \includegraphics[width=0.65\textwidth]{fig_manuscript_ci}
    \caption{Bootstrap 95\% confidence intervals for three models on the
    test set (n=100, 48 positives). The significance annotation refers to
    McNemar's test between multi-task and BiomedBERT v7 ($p=0.016$).}
    \label{fig:ci}
\end{figure}

\subsection{Threshold Analysis}

\begin{figure}[H]
    \centering
    \includegraphics[width=0.78\textwidth]{fig_threshold_curve}
    \caption{Precision, recall, and F1 versus decision threshold.
    The F1-optimal threshold ($\tau=0.13$) and a precision-oriented
    threshold ($\tau=0.50$) are annotated.}
    \label{fig:threshold}
\end{figure}

\subsection{Distribution Mismatch}

\begin{figure}[H]
    \centering
    \includegraphics[width=0.92\textwidth]{fig_distribution_mismatch}
    \caption{Distribution of interaction types in training positives (left,
    n=4,065 Qwen-labeled sentences) versus test set positives (right, n=48).}
    \label{fig:distribution}
\end{figure}

% ─────────────────────────────────────────────────────────────────────────────
\section{Discussion}
\label{sec:discussion}

\subsection{Soft Labels as the Primary Driver of Improvement}

\subsection{NER as Zero-Cost Auxiliary Supervision}

\subsection{Utility for BiotXplorer and Knowledge Graph Construction}

\subsection{The PubMed Distribution Ceiling}

\subsection{Limitations}

\begin{enumerate}
    \item \textbf{Test set size.}
    \item \textbf{Binary classification.}
    \item \textbf{English only.}
    \item \textbf{Automated NER labels.}
\end{enumerate}

% ─────────────────────────────────────────────────────────────────────────────
\section{Conclusion}

We have presented a multi-task BiomedBERT classifier for sentence-level
biotic interaction detection that achieves F1\,=\,0.868 on an ecological
test set. The model is built on three components: a two-stage teacher chain
that combines large-LLM binary annotation with ensemble soft probability
distillation; a multi-task NER auxiliary task that primes the shared encoder
at zero additional annotation cost; and a two-phase training protocol that
separates entity-level supervision from joint optimisation.

The central quantitative finding is that soft labels are the single most
impactful intervention in the pipeline, contributing +0.052 F1 on identical
architecture and data. The current performance ceiling at F1\,$\approx$\,0.87
is explained by a structural distribution mismatch between PubMed-sourced
training data and the ecological interaction types dominant in the test set,
pointing directly to the data collection strategy needed to break this ceiling.

The classifier is deployed in BiotXplorer
(\url{https://biotxplorer.sibils.org/}) as an upstream quality filter
and constitutes the first component of a knowledge graph construction
pipeline for biodiversity and host-pathogen interactions. Future work
will address: (1) harvesting from non-PubMed ecological literature to
expand training coverage of parasitism and host-interaction types;
(2) typed relation extraction to produce directed knowledge graph triples;
(3) integration with amplicon sequencing species detection for automated
interaction retrieval from biodiversity field surveys.

% ─────────────────────────────────────────────────────────────────────────────
\section*{Data and Code Availability}

Training code, evaluation scripts, and the test set are available at
\url{https://github.com/ecsltae/biotic-interaction-classifier}.
Model weights can be reproduced from the provided training configuration.

\bibliographystyle{abbrvnat}
\bibliography{references}

\end{document}
